\section{Danh sách kiểm tra khả năng tái lập}
\label{app:reproducibility}

\subsection{Các thành phần triển khai}

Các bước triển khai tương ứng với kết quả trong bài gồm:

\begin{itemize}
    \item \path{1_build_legal_causal_graph.py}: xây dựng đồ thị pháp luật;
    \item \path{2_build_causal_memory.py}: xây dựng memory cho rule/event và
    chỉ mục FAISS;
    \item \path{3_multi_hop_causal_retriever.py}: truy hồi seed, mở rộng và
    xếp hạng causal path;
    \item \path{causal_core/schema.py} và
    \path{causal_core/legal_scm.py}: canonical rule và structural inference;
    \item \path{4_counterfactual_verification.py}: query-aware verification,
    structural intervention và path ablation;
    \item \path{5_generate_final_answer.py}: chọn evidence và tạo câu trả lời;
    \item \path{5_5_generate_pipeline_predictions.py}: batch runner cho
    Steps~3--5;
    \item \path{7_compute_evaluation_metrics.py}: tính metric và diagnostic
    theo nhóm.
\end{itemize}

\subsection{Phiên bản và cấu hình chính}

\begin{table}[h]
    \centering
    \small
    \caption{Phiên bản và cấu hình cố định của thực nghiệm chính.}
    \label{tab:repro-settings}
    \begin{tabular}{@{}ll@{}}
        \toprule
        \textbf{Thành phần} & \textbf{Thiết lập} \\
        \midrule
        Benchmark & \texttt{2.0}, $250$ mẫu \\
        Step 4 & \texttt{4.1-query-aware-legal-scm} \\
        Step 5 & \texttt{5.1-query-aware-grounded-answer} \\
        Batch runner & \texttt{2.1-query-aware-answer} \\
        Evaluator & \texttt{2.0-current-pipeline-schema} \\
        Embedding & \texttt{BAAI/bge-m3} \\
        Retrieval hop / CF hop & $2$ / $3$ \\
        Answer provider & \texttt{extractive} \\
        Structural intervention & $\doop(M=\stateF)$ \\
        \bottomrule
    \end{tabular}
\end{table}

Các threshold đầy đủ được báo cáo trong
Bảng~\ref{tab:retrieval-hyperparameters},
Bảng~\ref{tab:verification-hyperparameters} và
Bảng~\ref{tab:generation-hyperparameters}.

\subsection{Input và paired artifact}

Các input chính gồm:

\begin{itemize}
    \item \path{data/blhs_multihop_benchmark_250.json};
    \item \path{data/blhs_rules_final_all_normalized.json};
    \item graph, causal memory và FAISS index được hai mode dùng chung.
\end{itemize}

Các prediction và evaluation directory được sử dụng để báo cáo kết quả là:

\begin{itemize}
    \item \path{data/pipeline_predictions_structural_scm_v41_step5_v51.json};
    \item \path{data/pipeline_predictions_path_ablation_v41_step5_v51.json};
    \item \path{evaluation_results/structural_scm_v41_step5_v51/};
    \item \path{evaluation_results/path_ablation_v41_step5_v51/}.
\end{itemize}

Artifact cũ tại
\path{evaluation_metrics_path_ablation_v41/} và prediction được tạo bởi
runner \texttt{2.0-primary-path-aware} không thuộc paired experiment chính
và không được dùng để tính các bảng trong bài.

\subsection{Mẫu lệnh chạy batch}

Với \texttt{MODE} lần lượt là \texttt{structural\_scm} và
\texttt{path\_ablation}, batch được chạy theo mẫu:

\begin{verbatim}
python 5_5_generate_pipeline_predictions.py ^
  --benchmark data/blhs_multihop_benchmark_250.json ^
  --provider extractive ^
  --counterfactual-mode MODE ^
  --rules data/blhs_rules_final_all_normalized.json ^
  --output data/pipeline_predictions_MODE_v41_step5_v51.json ^
  --jsonl-output data/pipeline_predictions_MODE_v41_step5_v51.jsonl ^
  --errors-output data/pipeline_errors_MODE_v41_step5_v51.json ^
  --run-log data/pipeline_run_log_MODE_v41_step5_v51.json ^
  --work-dir data/pipeline_intermediate_MODE_v41_step5_v51 ^
  --checkpoint-every 1 --resume
\end{verbatim}

Khi chạy trên Linux, ký tự nối dòng \texttt{\^{}} được thay bằng dấu gạch
chéo ngược. Step 7 được chạy với cùng benchmark, prediction tương ứng,
\texttt{--k-values 1 3 5 10} và \texttt{--strict}.

\subsection{Schema trace tối thiểu}

Mỗi prediction cần giữ tối thiểu các trường provenance sau:

\begin{itemize}
    \item sample ID, query và benchmark version;
    \item retrieval result, candidate path và primary path ID;
    \item query analysis, claim type và answer intent;
    \item final verification decision và verification score;
    \item factual/intervention state nếu dùng structural SCM;
    \item evidence ID, rule ID, event ID và article citation;
    \item Step 4, Step 5, runner và evaluator version;
    \item thời gian xử lý và lỗi, nếu có.
\end{itemize}

\subsection{Kiểm tra trước khi công bố}

Các mục sau vẫn phải được bổ sung từ đúng server và repository đã tạo
artifact; không được suy đoán:

\begin{itemize}
    \item hệ điều hành, phiên bản Python, CPU, RAM và GPU;
    \item revision chính xác của embedding model và dependency lockfile;
    \item immutable Git commit và repository URL công khai;
    \item hash của benchmark, rule corpus, graph, memory và prediction;
    \item random seed và trạng thái cache;
    \item giấy phép của code, dữ liệu và model;
    \item thông tin tác giả liên hệ, tài trợ và venue.
\end{itemize}

\todo{Điền thông tin phần cứng/server, commit, artifact hash và email tác giả
trước khi nộp bản thảo.}